\section{Every reconstruction row in explicit form}
\label{sec:inverse}

The main theorem of this section answers research question R2 of
\cite{incoming-bernoulli} for the complete $(2,2b)$ family. The answer
contains only factorials, Bernoulli numbers, and finite sums. In
particular, it does not leave an inverse matrix or an unknown
determinant in the definition of a reconstruction coefficient.

\subsection{A finite Pad\'e identity}

For $N\geq0$, define
\begin{equation}\label{eq:PN}
 P_N(t)=\sum_{k=0}^N
  \frac{(2N-k)!\,N!}{(2N)!\,k!\,(N-k)!}t^k.
\end{equation}
The exact initial Taylor coefficients of
$e^tP_N(-t)-P_N(t)$ vanish through degree $2N$, and
\begin{equation}\label{eq:Pade-remainder}
 e^tP_N(-t)-P_N(t)
 =\frac{(-1)^N(N!)^2}{(2N)!(2N+1)!}t^{2N+1}
   +O(t^{2N+2}).
\end{equation}
For completeness, the vanishing is the binomial identity
\[
 \sum_{k=0}^N(-1)^k\binom Nk(1+x)^{2N-k}
 =x^N(1+x)^N
\]
applied coefficient by coefficient. The first remaining coefficient is
\[
 \frac1{(2N)!}\sum_{k=0}^N
  \frac{(-1)^k\binom Nk}{2N+1-k}
 =\frac1{(2N)!}\int_0^1x^N(x-1)^N\dd x,
\]
which proves \eqref{eq:Pade-remainder} by the beta integral.

Set
\begin{equation}\label{eq:Fmu}
 F(z)=\sum_{j\geq0}\mu_jz^j
 =\frac{\sqrt z}{2}\coth\frac{\sqrt z}{2}.
\end{equation}
The right side is an even germ in $\sqrt z$, hence a single-valued
holomorphic germ in $z$. With $N=2n+1$, put
\[
 q_n(t^2)=\frac{P_{2n+1}(t)-P_{2n+1}(-t)}t,
 \qquad
 r_n(t^2)=\frac{P_{2n+1}(t)+P_{2n+1}(-t)}2.
\]
Then \eqref{eq:Pade-remainder} implies
\begin{equation}\label{eq:Fq-r}
 F(z)q_n(z)-r_n(z)=h_nz^{2n+1}+O(z^{2n+2}),
 \qquad
 h_n=\frac{((2n+1)!)^2}{(4n+2)!(4n+3)!}.
\end{equation}

\subsection{Two explicit orthogonal families}

Let $\mathcal L$ be the rational linear functional on $\Q[x]$ with
$\mathcal L(x^j)=\mu_j$. Define
\begin{align}
 \pi_n(x)&=\sum_{t=0}^n p_{nt}x^t,
 &p_{nt}&=
 \frac{2(2n+2t+1)!(2n+1)!}
 {(4n+2)!(2n-2t+1)!(2t)!},
 \label{eq:pi-coeff}\\
 c_n&=\frac{(2n+1)(2n+2)!}{(4n+2)!}.
 \label{eq:cn}
\end{align}
Thus $\pi_n(x)=x^nq_n(1/x)$ is monic. Coefficient comparison in
\eqref{eq:Fq-r} gives
\begin{equation}\label{eq:pi-moments}
 \mathcal L(\pi_n)=c_n,\qquad
 \mathcal L(x^k\pi_n)=0\ (1\leq k\leq n),\qquad
 \mathcal L(x^{n+1}\pi_n)=h_n.
\end{equation}
These are the inherited shifted-functional polynomials. The new inverse
construction requires orthogonality for $\mathcal L$ itself.

\begin{lemma}[The unshifted Bernoulli family]
\label{lem:Qorthogonal}
Put
\begin{align}
 Q_0(x)&=1,&\nu_0&=1,\notag\\
 Q_n(x)&=\pi_n(x)-\frac{c_n}{c_{n-1}}\pi_{n-1}(x),
 &\nu_n&=-\frac{c_n}{c_{n-1}}h_{n-1}\quad(n\geq1).
 \label{eq:Qnu}
\end{align}
Then the $Q_n$ are monic, pairwise orthogonal for $\mathcal L$, and
$\mathcal L(Q_n^2)=\nu_n\ne0$. Write
$q_{nt}=[x^t]Q_n(x)$, with $q_{nt}=0$ if $t>n$.
\end{lemma}
\begin{proof}
The zeroth moment of $Q_n$ vanishes by \eqref{eq:pi-moments}.
For $1\leq k<n$, both terms in
$\mathcal L(x^kQ_n)$ vanish by the same formula. Therefore $Q_n$
is orthogonal to every polynomial of degree less than $n$.
As $Q_n$ is monic,
\[
 \mathcal L(Q_n^2)=\mathcal L(x^nQ_n)
 =-\frac{c_n}{c_{n-1}}\mathcal L(x^n\pi_{n-1})
 =-\frac{c_n}{c_{n-1}}h_{n-1}.
\]
All factorials in $c_n,h_n$ are positive, so the norm is nonzero.
The functional is not positive: $\nu_0>0$ and $\nu_n<0$ for
$n\geq1$. The proof is algebraic and uses no positivity assumption.
\end{proof}

\begin{theorem}[Explicit inverse Bernoulli kernel]
\label{thm:kernel}
For $b\geq1$, let $H_b=(\mu_{t+k})_{0\leq t,k<b}$ and set
\begin{equation}\label{eq:kernel}
 K_b(x,y)=\sum_{n=0}^{b-1}\frac{Q_n(x)Q_n(y)}{\nu_n},
 \qquad
 K_{kt}^{(b)}=[x^ky^t]K_b(x,y)
 =\sum_{n=0}^{b-1}\frac{q_{nk}q_{nt}}{\nu_n}.
\end{equation}
Then $H_b$ is nonsingular and
$H_b^{-1}=(K_{kt}^{(b)})_{0\leq k,t<b}$.
Moreover,
\begin{equation}\label{eq:Kfirst}
 K_{0t}^{(b)}=\frac{p_{b-1,t}}{c_{b-1}}.
\end{equation}
\end{theorem}
\begin{proof}
Every polynomial $f$ of degree less than $b$ has the orthogonal
expansion
\[
 f(x)=\sum_{n=0}^{b-1}
   \frac{\mathcal L(fQ_n)}{\nu_n}Q_n(x).
\]
It follows that $\mathcal L_y(K_b(x,y)f(y))=f(x)$.
Taking $f(y)=y^j$ and comparing coefficients of $x^k$ gives
$\sum_tK_{kt}^{(b)}\mu_{t+j}=\delta_{kj}$, which proves the matrix
identity. The polynomial $\pi_{b-1}/c_{b-1}$ has moments one at
degree zero and zero at degrees $1,\ldots,b-1$ by
\eqref{eq:pi-moments}; uniqueness of the first inverse row proves
\eqref{eq:Kfirst}.
\end{proof}

\subsection{The complete mixed inverse}

Define the following finite rational data, for $0\leq t,k<b$:
\begin{align}
 u_t^{(b)}&=\mu_t
   \frac{(2b)!(2t+1)!}{2(2b+2t-1)!},
 \label{eq:ut}\\
 v_k^{(b)}&=\sum_{t=0}^{b-1}K_{kt}^{(b)}u_t^{(b)},
 &\delta_b&=2-B_{2b-2}.
 \label{eq:vkdelta}
\end{align}
Here every $K_{kt}^{(b)}$ is already given by the explicit finite sum
\eqref{eq:kernel} and the factorial coefficients
\eqref{eq:pi-coeff}--\eqref{eq:Qnu}.

\begin{theorem}[All reconstruction rows]
\label{thm:all-rows}
For every $b\geq2$, $\delta_b\ne0$, and
\begin{equation}\label{eq:all-rows}
 \boxed{\displaystyle
 \mathcal R_{b;k,t}=
 \frac{(2b-2k)!}{(2b+2t-1)!}
 \left(
 K_{kt}^{(b)}-
 \frac{v_k^{(b)}}{\delta_b}\frac{p_{b-1,t}}{c_{b-1}}
 \right)
 =(M_b^{-1})_{kt}.}
\end{equation}
Consequently every even-zeta coordinate is recovered by
\begin{equation}\label{eq:all-coordinates}
 \boxed{\displaystyle
 \zeta(2k)=\sum_{t=0}^{b-1}\mathcal R_{b;k,t}C_{b+t},
 \qquad 0\leq k<b.}
\end{equation}
This gives all rows in a common closed finite-sum formula.
\end{theorem}
\begin{proof}
Let
\[
 R_b=\operatorname{diag}_{t}(2b+2t-1)!,\qquad
 E_b=\operatorname{diag}_{k}\frac1{(2b-2k)!}.
\]
The exact entries in \eqref{eq:M-system} give
\begin{equation}\label{eq:rank-one}
 M_b=R_b(H_b+u^{(b)}e_0^{\mathsf T})E_b.
\end{equation}
By \eqref{eq:Kfirst}, factorial cancellation gives
\begin{align}\label{eq:v0}
 v_0^{(b)}
 &=\frac1{2b-1}\sum_{t=0}^{b-1}
   \binom{2b-1}{2t}(2t+1)B_{2t}
 =1-B_{2b-2}.
\end{align}
To prove the last equality directly, put $N=2b-1\geq3$ and use
\[
 \sum_{j=0}^N\binom NjB_j=B_N(1)=0,\qquad
 \sum_{j=0}^N j\binom NjB_j=-NB_{N-1}.
\]
After removing the $j=1$ terms and using the vanishing of the higher
odd Bernoulli numbers, their sum is exactly \eqref{eq:v0}.

We include the elementary nonvanishing argument. In the localization
$\Z_{(2)}$ of rationals with odd denominator, the Bernoulli recurrence
is
\[
 (2m+1)(2B_{2m})=2m-1-
  \sum_{j=1}^{m-1}\binom{2m+1}{2j}(2B_{2j}).
\]
Induction and
$\sum_{j=1}^{m-1}\binom{2m+1}{2j}=2^{2m}-2m-2$
show that $2B_{2m}$ is an odd element of $\Z_{(2)}$.
It cannot equal four. Hence $B_{2m}\ne2$, proving $\delta_b\ne0$.

Let $K=H_b^{-1}$ and $v=Ku$. A direct multiplication, using
$1+e_0^{\mathsf T}v=\delta_b$, shows
\[
 (H_b+ue_0^{\mathsf T})^{-1}
 =K-\frac{ve_0^{\mathsf T}K}{\delta_b}.
\]
Insert \eqref{eq:Kfirst} and multiply by the two inverse diagonal
factors from \eqref{eq:rank-one}. This is \eqref{eq:all-rows}.
Finally apply it to \eqref{eq:M-system}.
\end{proof}

\begin{remark}[Recovery of the earlier rational row]
The row $k=0$ simplifies to
\begin{equation}\label{eq:rational-row}
 \mathcal R_{b;0,t}=
 \frac{2}{(2b-1)(2-B_{2b-2})}\binom{2b-1}{2t}.
\end{equation}
Thus the previously known rational extraction is recovered exactly.
Rows $k=1,\ldots,b-1$ in \eqref{eq:all-rows} are the missing rows
asked for in the source. The first nontrivial complete inverse is
\begin{equation}\label{eq:b2inverse}
 M_2=
 \begin{pmatrix}3/4&1/4\\2/3&-1/12\end{pmatrix},\qquad
 M_2^{-1}=\frac1{11}
 \begin{pmatrix}4&12\\32&-36\end{pmatrix}.
\end{equation}
These small matrices are included as a transparent normalization check.
\end{remark}
